\zag{Pojęcie i rodzaje procesu karnego}[{\kolor A}]\str{25–29,\\{\kolor tu przerwa odpowiedniej długosci} 60–65}
% \textit{To dość obszerne}\\\medskip
Proces karny jest to (w znaczeniu \textbf{ogólnym}) uregulowany prawem ciąg czynności jego uczestników mający na celu rozpoznanie sprawy karnej, a więc wykrycie i rozpoznanie przestępstwa i jego sprawcy, wyjaśnienie okoliczności i ustalenie konsekwencji prawnych oraz pociągnięcie winnego do odpowiedzialności.\medskip

Proces karny w znaczeniu \textbf{konkretnym} jest działalnością w konkretnej sprawie (proces jakiegoś Xińskiego) i powinien przebiegać w zgodzie ze wzorcem jakim jest proces karny w znaczeniu ogólnym.\medskip

Ogół procesów w znaczeniu konkretnym składa się na proces karny w znaczeniu \textbf{rzeczywistym}, jest to więc ogólna praktyka procesu karnego. Rozbierzności między nim a procesem w znaczeniu ogólnym wskazują więc na ogólne i powtarzające się odstępstwa uczestników procesów w znaczeniu konkretnym od normatywnego modelu wyrażonego w ramach procesu w znaczeniu ogólnym.\medskip

Proces karny może dotyczyć zarówno odpowiedzialności karnej oskarżonego jak i jego odpowiedzialności cywilnej. Postępowanie dotyczące odpowiedzialności karnej to \textbf{proces zasadniczy}, to w sprawie odpowiedzialności cywilnej to \textbf{akcja cywilna}.\medskip

Ze względu na sąd właściwy do rozpoznania sprawy rozróżnić można \textbf{proces powszechny} (przed sądami powszechnymi) od \textbf{procesu specjalnego} (przed sądami wojskowymi oraz Trybunałem Stanu).\medskip

Innym sposobem różnicowania jest zastosowanie \textbf{kryterium trybu oskarżenia} (publicznego lub prywatnego) albo trybu samego postępowania (\textbf{tryb zwyczajny} i \textbf{tryb szczególny}). Tryb zwyczajny definiuje się \textit{idem per idem}, natomiast trybem szczególnym jest ten, który od normy tej odbiega.\medskip

Same tryby szczególne dzielić można według różnych kryteriów, np: obowiązkowość (fakultatywny a obligatoryjny), \textit{rationae materiae} (w sprawach wielkiej wagi a w sprawach drobnych), \textit{rationae personae} (nieletni, wojskowi), sąd orzekający (sądy powszechne a sądy szczególne), źródło prawa (\acr{kpk} a inne ustawy) lub wpływ na formalizm (postępowania ekwiwalentne, a więc tak samo sformalizowane jak tryb zwyczajny, postępowania wzbogacone, czyli bardziej sformalizowane i postępowania zredukowane, czyli o zmniejszonym formalizmie).


% TEGO NIE:
% istota i aksjologia
% cele
% przedmiot
% pojeęcie i warunki rzetelności
% postępowanie karne, jego rodzaje, mediacyjne
% granice
